\chapter{Four-Dimensional N=2 Gauge Dynamics and Seiberg-Witten Theory}
\label{ch:susy-four-dimensional-n2-seiberg-witten}

\ConventionsLorentzian

This chapter develops the first layer of four-dimensional \(\mathcal N=2\)
gauge dynamics.  The emphasis is the Wilsonian low-energy description on the
Coulomb branch, where the exact two-derivative Abelian action is encoded by a
holomorphic prepotential and by periods of a family of curves.

\section{Coulomb-Branch Coordinates}

Let \(G\) be a compact connected gauge group with Lie algebra \(\mathfrak g\).
An \(\mathcal N=2\) vector multiplet contains a gauge field, two gauginos, and
a complex scalar \(\phi\in\mathfrak g_{\C}\).  A Coulomb-branch point is a
vacuum for which \(\phi\) has expectation value in a Cartan subalgebra and
the low-energy gauge group contains an Abelian factor \(U(1)^r\), where
\(r=\operatorname{rank}G\).  Gauge-invariant holomorphic coordinates are
polynomial functions of \(\phi\), for example
\[
  u_k=\left\langle \operatorname{tr}(\phi^k)\right\rangle .
\]
The trace convention is the convention fixed in the gauge chapters: the
Yang--Mills kinetic term uses
\[
  \frac{1}{4g^2}\operatorname{tr}(F_{\mu\nu}F^{\mu\nu})
\]
with the chosen invariant bilinear form.

\section{Special Coordinates And Prepotential}

On a smooth patch of the Coulomb branch, choose electric special coordinates
\[
  a^I,\qquad I=1,\ldots,r .
\]
The two-derivative Wilsonian action is determined by a holomorphic function
\(\mathcal F(a)\), the prepotential, through
\[
  a_{D,I}=\frac{\partial\mathcal F}{\partial a^I},
  \qquad
  \tau_{IJ}(a)=\frac{\partial^2\mathcal F}{\partial a^I\partial a^J}.
\]
The imaginary part of \(\tau_{IJ}\) is the positive gauge-coupling matrix in
the region where the Abelian Wilsonian description is valid.  Electric-magnetic
duality acts on the column vector
\[
  \Pi=\begin{pmatrix} a_D \\ a \end{pmatrix}
\]
by integral symplectic matrices.  If \(\gamma=(n_m^I,n_{e,I})\) is an
electromagnetic charge vector, its central charge is
\[
  Z_\gamma=n_{e,I}a^I+n_m^I a_{D,I}+\sum_\alpha s_\alpha m_\alpha ,
\]
where \(m_\alpha\) are background flavor masses and \(s_\alpha\) are flavor
charges.  The BPS inequality has the convention
\[
  M_\gamma\geq \sqrt2\,|Z_\gamma|.
\]
States saturating this bound form shortened \(\mathcal N=2\) multiplets.

\section{Seiberg-Witten Period Data}

A Seiberg-Witten period datum consists of a family of compact Riemann surfaces
\(\Sigma_u\), a meromorphic differential \(\lambda_{\mathrm{SW}}\), and a
symplectic basis of one-cycles \((A_I,B^I)\) such that
\[
  a^I(u)=\frac{1}{2\pi}\oint_{A_I}\lambda_{\mathrm{SW}},
  \qquad
  a_{D,I}(u)=\frac{1}{2\pi}\oint_{B^I}\lambda_{\mathrm{SW}} .
\]
The differential is fixed up to exact terms by requiring that
\(\partial\lambda_{\mathrm{SW}}/\partial u_k\) span holomorphic differentials
on \(\Sigma_u\) and that the large-field asymptotics reproduce the
one-loop Wilsonian coupling.

\begin{proposition}[Integrability of the period coupling]
\label{prop:sw-period-integrability}
Assume that the period matrix \(\tau_{IJ}=\partial a_{D,I}/\partial a^J\)
comes from a family of curves with symplectic intersection pairing and varies
holomorphically on a simply connected patch where the \(a^I\) are coordinates.
Then locally there is a holomorphic prepotential \(\mathcal F(a)\) satisfying
\(a_{D,I}=\partial\mathcal F/\partial a^I\).
\end{proposition}

\begin{proof}
The Riemann bilinear relation implies symmetry
\(\partial a_{D,I}/\partial a^J=\partial a_{D,J}/\partial a^I\).  Therefore
the holomorphic one-form \(\sum_I a_{D,I}\,da^I\) on the coordinate patch is
closed.  On a simply connected patch every closed holomorphic one-form is
exact, so there exists \(\mathcal F\) with
\(d\mathcal F=\sum_I a_{D,I}\,da^I\).  This is the required prepotential.
\end{proof}

\section{Pure Gauge Theory With Gauge Algebra su two}

For the pure theory with gauge algebra \(\mathfrak{su}(2)\), the Coulomb
coordinate is
\[
  u=\left\langle \operatorname{tr}\phi^2\right\rangle .
\]
The Weyl group identifies \(a\sim -a\), and the large-field relation is
\[
  u=\frac12a^2+O(\Lambda^4/a^2),
  \qquad
  a(u)\sim\sqrt{2u}.
  \label{eq:sw-large-u-a-u}
\]
The perturbative Wilsonian prepotential in the pure theory has the one-loop
form
\begin{equation}
  \mathcal F_{\rm pert}(a)
  =
  \frac{\ii}{2\pi}a^2
  \log\frac{a^2}{\Lambda^2}
  +
  \hbox{terms polynomial in }a .
  \label{eq:sw-perturbative-prepotential}
\end{equation}
The polynomial terms change \(a_D\) by an integral affine electric
redefinition and do not change the monodromy class.  Thus
\[
  a_D=\frac{\partial\mathcal F}{\partial a}
  =
  \frac{\ii}{\pi}a
  \log\frac{a^2}{\Lambda^2}
  +O(a).
  \label{eq:sw-large-u-ad}
\]
When \(u\) circles infinity once counterclockwise, the square root in
\eqref{eq:sw-large-u-a-u} gives \(a\mapsto -a\), while the logarithm in
\eqref{eq:sw-large-u-ad} shifts by \(2\pi\ii\).  Therefore the period vector
\(\Pi=(a_D,a)^T\) has monodromy
\begin{equation}
  M_\infty
  =
  \begin{pmatrix}
  -1&2\\
  0&-1
  \end{pmatrix}.
  \label{eq:sw-infinity-monodromy-derived}
\end{equation}

The anomalous \(U(1)_R\) symmetry of the microscopic theory is reduced to
\(\mathbb Z_8\).  A point on the Coulomb branch with coordinate \(u\) further
breaks it to \(\mathbb Z_4\), leaving the quotient action \(u\mapsto -u\).
Consequently finite singularities, if there are only two in the minimal
period problem, occur at opposite values.  The scale \(\Lambda\) fixes their
locations to \(u=\pm\Lambda^2\).  At a nodal singular fiber, a primitive
charge \(\gamma=(n_m,n_e)\) becomes massless and the Picard-Lefschetz
transformation of the period vector is
\begin{equation}
  M_\gamma
  =
  \begin{pmatrix}
  1+2n_m n_e&2n_e^2\\
  -2n_m^2&1-2n_m n_e
  \end{pmatrix}.
  \label{eq:sw-picard-lefschetz-charge-monodromy}
\end{equation}
The factor \(2\) is the hypermultiplet contribution in the charge convention
where the central charge is \(Z_\gamma=n_ea+n_m a_D\).  Choosing the
singularity at \(u=\Lambda^2\) to be a monopole point, \(\gamma_m=(1,0)\),
gives
\[
  M_{\Lambda^2}=M_{(1,0)}
  =
  \begin{pmatrix}
  1&0\\
  -2&1
  \end{pmatrix}.
\]
The second finite monodromy is fixed by
\[
  M_{\Lambda^2}M_{-\Lambda^2}=M_\infty,
\]
hence
\begin{equation}
  M_{-\Lambda^2}
  =
  \begin{pmatrix}
  -1&2\\
  -2&3
  \end{pmatrix}
  =
  M_{(1,-1)} .
  \label{eq:sw-dyon-monodromy-derived}
\end{equation}
Thus the second singularity carries a massless dyon of charge
\(\gamma_d=(1,-1)\) in this period basis.

The elliptic curve must have a discriminant with precisely these two finite
zeros and must reproduce the weak-coupling cusp at \(u=\infty\).  In a
hyperelliptic genus-one coordinate \(x\), after rescaling \(x\) and \(u\), the
minimal such family is
\begin{equation}
  \Sigma_u:\qquad
  y^2=(x^2-\Lambda^4)(x-u).
  \label{eq:sw-su2-curve-derived}
\end{equation}
Its branch points are
\[
  x=\Lambda^2,\qquad x=-\Lambda^2,\qquad x=u,\qquad x=\infty,
\]
and the finite discriminant is proportional to
\[
  \Delta(u)=u^2-\Lambda^4 .
\]
At \(u=\Lambda^2\) the branch points \(x=u\) and \(x=\Lambda^2\) collide; at
\(u=-\Lambda^2\) the branch points \(x=u\) and \(x=-\Lambda^2\) collide.  The
monodromies of the corresponding vanishing cycles are exactly
\eqref{eq:sw-picard-lefschetz-charge-monodromy} with the charges above.  This
derivation uses the following hypotheses: the rank-one Coulomb branch has
only the two finite singularities required by the residual \(R\)-symmetry, the
singular fibers are nodal hypermultiplet singularities, and the weak-coupling
monodromy is exhausted by the perturbative vector multiplet.

\section{Periods, Picard-Fuchs Equation, and BPS Masses}

Choose the branch of the periods in the weak-coupling region
\(|u|\gg|\Lambda|^2\) by
\[
  a(u)\sim\sqrt{2u}.
\]
With the curve \eqref{eq:sw-su2-curve-derived}, the periods obey the
Picard-Fuchs equation
\begin{equation}
  (u^2-\Lambda^4)\frac{d^2\Pi}{du^2}
  +
  \frac14\Pi=0,
  \qquad
  \Pi\in\{a_D,a\}.
  \label{eq:sw-picard-fuchs-equation}
\end{equation}
The solution normalized as the electric period at infinity can be written as
\begin{equation}
  a(u)
  =
  \sqrt{2(u+\Lambda^2)}\,
  {}_2F_1
  \left(
    -\frac12,\frac12;1;
    \frac{2\Lambda^2}{u+\Lambda^2}
  \right),
  \label{eq:sw-electric-period-hypergeometric}
\end{equation}
with the branch chosen by \eqref{eq:sw-large-u-a-u}.  A dual period that
vanishes at the monopole point is
\begin{equation}
  a_D^{(m)}(u)
  =
  \frac{\ii}{2\Lambda}(u-\Lambda^2)
  {}_2F_1
  \left(
    \frac12,\frac12;2;
    \frac{\Lambda^2-u}{2\Lambda^2}
  \right).
  \label{eq:sw-monopole-period-hypergeometric}
\end{equation}
The superscript records the local branch near \(u=\Lambda^2\).  Analytic
continuation of the pair
\((a_D^{(m)},a)\) around the singularities gives the monodromy matrices above.
Near the monopole point,
\begin{equation}
  a_D^{(m)}(u)
  =
  \frac{\ii}{2\Lambda}(u-\Lambda^2)
  +O((u-\Lambda^2)^2).
  \label{eq:sw-monopole-period-linear}
\end{equation}
The monopole central charge is \(Z_{(1,0)}=a_D\), hence
\begin{equation}
  M_{(1,0)}(u)
  =
  \sqrt2\,|a_D(u)|
  =
  \frac{|u-\Lambda^2|}{\sqrt2|\Lambda|}
  +O(|u-\Lambda^2|^2).
  \label{eq:sw-monopole-mass-near-singularity}
\end{equation}
At \(u=-\Lambda^2\), use the local period basis obtained by transporting to
that singularity.  The vanishing central charge is
\[
  Z_{(1,-1)}=a_D-a,
\]
and the same nodal-fiber analysis gives
\begin{equation}
  a_D(u)-a(u)
  =
  c_-\,(u+\Lambda^2)
  +
  O((u+\Lambda^2)^2),
  \qquad
  c_-\neq0,
  \label{eq:sw-dyon-period-linear}
\end{equation}
on the dyon branch.  Therefore
\begin{equation}
  M_{(1,-1)}(u)
  =
  \sqrt2\,|a_D(u)-a(u)|
  =
  \sqrt2\,|c_-|\,|u+\Lambda^2|
  +O(|u+\Lambda^2|^2).
  \label{eq:sw-dyon-mass-near-singularity}
\end{equation}
The phase and magnitude of \(c_-\) depend on the branch and symplectic basis;
the linear vanishing and the charge \((1,-1)\) are invariant data of the
singular fiber.

\section{ADHM Moduli And The Nekrasov Expansion}
\label{sec:n2-adhm-nekrasov-expansion}

The Seiberg-Witten curve determines the low-energy prepotential indirectly
through periods.  There is also a weak-coupling construction of the same
instanton expansion from the ultraviolet gauge theory.  This section records
the construction with all convention-dependent inputs exposed.

Work in Euclidean signature on \(\mathbb R^4\simeq\mathbb C^2\), with the
orientation and instanton number convention of
Section~\ref{sec:bpst-instanton-thooft-vertex}.  The integer
\[
  k=\frac1{8\pi^2}\int_{\mathbb R^4}\operatorname{tr}(F\wedge F)
\]
is positive for the orientation chosen there.  The instanton-counting
coordinate will be denoted by \(q\).  In the common half-trace superspace
coordinate, \(q=\exp(2\pi\ii\tau_{\rm ht})\).  In the trace-delta component
coupling coordinate of Section~\ref{sec:bpst-instanton-thooft-vertex}, the
same one-instanton exponential is
\[
  q
  =
  \exp\!\left(
    -\frac{4\pi^2}{g_{\rm YM}^2}+\ii\theta
  \right)
  =
  \exp\!\left(
    -\frac{8\pi^2}{g_{\rm ht}^2}+\ii\theta
  \right).
\]
The ADHM and localization formulae below use \(q\) only as the holomorphic
coordinate that grades instanton number.  For the pure \(SU(N)\) theory,
dimensional transmutation replaces \(q\) by \(\Lambda^{2N}\) after the
renormalization scheme has been fixed.

\begin{definition}[Framed ADHM data for \(U(N)\)]
\label{def:framed-adhm-data-un}
Let \(K\simeq\mathbb C^k\) and \(W\simeq\mathbb C^N\).  Framed ADHM data are
quadruples
\[
  (B_1,B_2,I,J),
  \qquad
  B_1,B_2\in\operatorname{End}(K),\quad
  I\in\operatorname{Hom}(W,K),\quad
  J\in\operatorname{Hom}(K,W),
\]
obeying
\begin{align}
  [B_1,B_2]+IJ&=0,
  \label{eq:adhm-complex-moment-map}\\
  [B_1,B_1^\dagger]+[B_2,B_2^\dagger]+II^\dagger-J^\dagger J
  &=\zeta\,\mathbf 1_K .
  \label{eq:adhm-real-moment-map}
\end{align}
The quotient by \(U(K)\) is the hyper-K\"ahler quotient.  The parameter
\(\zeta>0\) is a stability or noncommutative resolution parameter; the
ordinary framed instanton moduli space is recovered at \(\zeta=0\) away from
small-instanton singular strata.  The framing is a trivialization at
\(S^3_\infty\).
\end{definition}

\begin{proposition}[ADHM dimension count]
\label{prop:adhm-dimension-count}
On the stable locus where the \(GL(K)\) action has no stabilizer, the framed
\(k\)-instanton ADHM moduli space for \(U(N)\) has complex dimension
\[
  2kN
\]
and real dimension \(4kN\).  Removing the center-of-mass position in
\(\mathbb R^4\) gives the centered dimension \(4kN-4\).
\end{proposition}

\begin{proof}
The complex variables \(B_1,B_2,I,J\) have complex dimension
\[
  2k^2+kN+kN=2k^2+2kN .
\]
The complex moment-map equation
\eqref{eq:adhm-complex-moment-map} is an equation in
\(\operatorname{End}(K)\), hence imposes \(k^2\) complex equations on the
smooth stable locus.  Quotienting by the complexified gauge group
\(GL(K)\) removes another \(k^2\) complex dimensions.  Thus
\[
  (2k^2+2kN)-k^2-k^2=2kN .
\]
The real moment-map quotient by \(U(K)\) is equivalent to this stable
complex quotient by the Kempf--Ness correspondence.  Therefore the real
dimension is \(2(2kN)=4kN\).  Translation of the instanton center is the
trace part of \(B_1,B_2\), a copy of \(\mathbb C^2\simeq\mathbb R^4\);
fixing it removes four real dimensions.
\end{proof}

\begin{remark}[Gauge group and compactification]
\label{rem:adhm-u-n-su-n-compactification}
The \(U(N)\) ADHM quotient is the clean equivariant model because the framing
torus has \(N\) weights \(a_1,\ldots,a_N\).  Pure \(SU(N)\) instanton
calculations are obtained by imposing \(\sum_{\alpha=1}^Na_\alpha=0\) and by
discarding the decoupled abelian center in observables.  Ordinary
commutative \(U(1)\) gauge theory has no smooth finite-action instantons on
\(\mathbb R^4\); the \(U(1)\) fixed points in the sheaf or
noncommutative compactification are pointlike ideal-sheaf strata.  Thus the
compactification or contour prescription is part of the definition of
\(Z_{\rm inst}\), as in Definition~\ref{def:s4-nekrasov-factor}.
\end{remark}

Let
\[
  T=(\mathbb C^\times)^2_{\epsilon_1,\epsilon_2}
  \times(\mathbb C^\times)^N_{a_1,\ldots,a_N}
\]
act by rotating \(\mathbb C^2\) and by the framing torus.  The parameters
\(\epsilon_1,\epsilon_2\) are the equivariant weights of the two complex
planes, and \(a_\alpha\) are Coulomb weights.  The equivariant deformation is
called the \(\Omega\)-background.  It has two jobs: it makes the instanton
volume integral an equivariant integral, and it isolates fixed points for
localization.

\begin{quotedtheorem}[ADHM fixed points]
\label{qt:adhm-fixed-points-young-diagrams}
For the smooth framed \(U(N)\) ADHM compactification on \(\mathbb C^2\), with
generic \(a_\alpha,\epsilon_1,\epsilon_2\), the \(T\)-fixed points are
isolated and are labeled by \(N\)-tuples of Young diagrams
\[
  \vec Y=(Y_1,\ldots,Y_N),
  \qquad
  |\vec Y|:=\sum_{\alpha=1}^N |Y_\alpha|=k .
\]
The equivariant integral over the \(k\)-instanton component equals the sum of
the reciprocal equivariant Euler classes of the tangent complexes at these
fixed points.
\end{quotedtheorem}

The quoted statement is a theorem of equivariant algebraic geometry for the
framed sheaf compactification.  Its use in the physical path integral
additionally assumes that the supersymmetric regulator, the
\(\Omega\)-background, and the contour define the same equivariant integral.
This physical identification is the localization input; the dimension count
alone does not imply it.

\begin{definition}[Pure vector Nekrasov partition function]
\label{def:pure-vector-nekrasov-partition}
For pure \(U(N)\) vector multiplets, define
\[
  Z_{\rm inst}^{U(N)}(a,\epsilon_1,\epsilon_2;q)
  =
  \sum_{k\geq0}q^k Z_k^{U(N)}(a,\epsilon_1,\epsilon_2),
  \qquad
  Z_0^{U(N)}=1,
\]
where \(Z_k^{U(N)}\) is the equivariant integral over the framed
\(k\)-instanton compactification.  The one-instanton fixed-point formula is
\begin{equation}
  Z_1^{U(N)}
  =
  \frac1{\epsilon_1\epsilon_2}
  \sum_{\alpha=1}^N
  \prod_{\beta\neq\alpha}
  \frac1{(a_\alpha-a_\beta)
  (a_\alpha-a_\beta+\epsilon_1+\epsilon_2)} .
  \label{eq:nekrasov-one-instanton-un}
\end{equation}
The pure \(SU(N)\) specialization imposes \(\sum_\alpha a_\alpha=0\).  Matter
hypermultiplets multiply each fixed-point term by the equivariant Euler class
of the corresponding Dirac complex; the vector formula above is the part
needed for the pure \(SU(2)\) Seiberg-Witten comparison.
\end{definition}

\begin{proposition}[Pure \(SU(2)\) one-instanton coefficient]
\label{prop:nekrasov-su2-one-instanton}
Set \(a_1=a\), \(a_2=-a\), and
\(\epsilon=\epsilon_1+\epsilon_2\).  Then
\[
  Z_1^{SU(2)}
  =
  \frac{2}{\epsilon_1\epsilon_2(4a^2-\epsilon^2)} .
\]
Consequently, in the Nekrasov normalization
\[
  \mathcal F_{\rm inst}(a;q)
  :=
  \lim_{\epsilon_1,\epsilon_2\to0}
  \epsilon_1\epsilon_2
  \log Z_{\rm inst}^{SU(2)}(a,\epsilon_1,\epsilon_2;q),
\]
the first pure-gauge term is
\[
  \mathcal F_{\rm inst}(a;q)
  =
  \frac{q}{2a^2}+O(q^2).
\]
For the pure \(SU(2)\) theory, \(q\) is identified with \(\Lambda^4\) in the
weak-coupling scale coordinate.
\end{proposition}

\begin{proof}
Substituting \(a_1=a\), \(a_2=-a\) in
\eqref{eq:nekrasov-one-instanton-un} gives two fixed-point contributions:
\[
  Z_1^{SU(2)}
  =
  \frac1{\epsilon_1\epsilon_2}
  \left[
    \frac1{(2a)(2a+\epsilon)}
    +
    \frac1{(-2a)(-2a+\epsilon)}
  \right].
\]
The second denominator is \(2a(2a-\epsilon)\).  Hence
\[
  Z_1^{SU(2)}
  =
  \frac1{\epsilon_1\epsilon_2}
  \frac1{2a}
  \left(
    \frac1{2a+\epsilon}
    +
    \frac1{2a-\epsilon}
  \right)
  =
  \frac{2}{\epsilon_1\epsilon_2(4a^2-\epsilon^2)} .
\]
Since \(Z_{\rm inst}=1+qZ_1+O(q^2)\),
\[
  \epsilon_1\epsilon_2\log Z_{\rm inst}
  =
  q\,\frac{2}{4a^2-\epsilon^2}
  +O(q^2).
\]
Taking \(\epsilon_1,\epsilon_2\to0\) gives \(q/(2a^2)\) at order \(q\).
\end{proof}

The coefficient in Proposition~\ref{prop:nekrasov-su2-one-instanton} is the
first weak-coupling instanton correction to the same holomorphic prepotential
whose periods are encoded by the curve \eqref{eq:sw-su2-curve-derived},
after matching the additive quadratic ambiguity in \(\mathcal F\) and the
scale coordinate.  The comparison is stringent because the Seiberg-Witten
curve determines all powers of \(\Lambda^4/a^4\) by a period problem, while
the ADHM construction determines them by finite-dimensional equivariant
integrals.

The exponential size of the \(k\)-instanton sector is \(q^k\).  In a
transseries analysis of the weak-coupling expansion, these actions mark the
candidate locations of Borel singularities separated from the perturbative
sector by integer multiples of the one-instanton action.  A theorem-level
resurgence statement for four-dimensional supersymmetric gauge theory would
require a regulator, a summation prescription, and uniform control of the
large-order perturbative coefficients.  The assertion here is the convention
ledger: the same \(k\), action, and \(\theta\)-phase grade the ADHM integral,
the \(\Omega\)-background partition function, and the instanton transseries
sectors.

\section{Monodromy And Light Charges}

The weak-coupling region is the exterior of the curve of marginal stability
on which \(a_D/a\) is real.  In the charge convention fixed above, the
semiclassical BPS spectrum contains the vector multiplet with charge
\((0,1)\) and a tower of hypermultiplets with charges
\begin{equation}
  (n_m,n_e)=(1,n),
  \qquad
  n\in\mathbb Z,
  \label{eq:sw-weak-coupling-bps-tower}
\end{equation}
together with antiparticles.  The tower is generated by monodromy at infinity
from either finite singular charge.  The central charges are
\[
  Z_{(0,1)}=a,
  \qquad
  Z_{(1,n)}=a_D+na .
\]
Inside the strong-coupling chamber bounded by the marginal-stability curve,
the stable hypermultiplets are the two light states associated with the two
singularities, with charges
\[
  (1,0),\qquad (1,-1),
\]
and their antiparticles.  The change in the BPS index occurs when
\[
  \arg Z_{\gamma_1}=\arg Z_{\gamma_2}
\]
for charges that can form or decay as a BPS bound state.  This is the
wall-crossing mechanism in its rank-one Seiberg-Witten form: the central
charge function is continuous, while the protected index of one-particle BPS
states is locally constant only away from marginal-stability walls.

\section{Low-Energy Abelian Theory Near A Singular Point}

Near a point where the monopole period \(a_D\) vanishes, the correct
low-energy variables are the dual \(U(1)\) vector multiplet with scalar
\(a_D\) and the light monopole hypermultiplet.  A local Wilsonian
superpotential has the form
\[
  W_{\mathrm{loc}}=\sqrt2\,A_D M\widetilde M ,
\]
where \(A_D\) is the dual chiral multiplet and \(M,\widetilde M\) are monopole
chiral multiplets of opposite dual electric charge.  Adding a small
\(\mathcal N=1\)-preserving deformation proportional to \(u\) gives an
F-term equation that forces monopole condensation in the local theory.  This
is the microscopic input behind the Abelian confinement mechanism used later
in the Witten--Donaldson and Seiberg-Witten comparison.

\section{Status Boundary}

The period construction determines an Abelian Wilsonian effective action under
explicit hypotheses: existence of the \(\mathcal N=2\) continuum theory,
existence of the Coulomb-branch vacuum family, validity of the two-derivative
Abelian description away from singular loci, and BPS saturation for the
states assigned to vanishing cycles.  The missing nonperturbative regulator
for four-dimensional supersymmetric gauge theory remains part of the
Wilsonian-scheme problem recorded in Chapter~\ref{ch:susy-four-dimensional-n1-gauge-dynamics}.

\begin{openproblem}[Theorem-level Seiberg-Witten construction]
\label{op:theorem-level-seiberg-witten-construction}
Construct the four-dimensional \(\mathcal N=2\) gauge theory with a regulator
and continuum limit sufficient to derive the Coulomb-branch special geometry,
period data, BPS spectrum near singular fibers, and low-energy Abelian
Wilsonian action as theorems of QFT.
\end{openproblem}
